% Nouveau EGA V (cote 25): archive cote 25, site page 90 (PDF page 91).
% typescript. Draft compared: none.
% AI-transcribed by Claude Opus 5.5 (claude-opus-5-5, adaptive thinking, effort high), Message Batches API, batch ega25b (msgbatch_01A7n5YuXubUPNK1zsepi8RG), 2026-11-13; not yet checked by eye.
\scanpage{25/90}
\mspage{\unsure{87}}

\uline{Proposition} 14.5. Supposons dim X $\leqslant$ m-1, \xout{Soit Y l'image fermée}
\xout{\illegible{1}} f:X $\to$ P une immersion, \xout{Soit} enfin X séparable sur k. Soit Y\hand{$_y$} l'image
schématique de $X_{k(\hand{y})}$ dans $Q(C_{\hand{y}})$. Alors le morphisme induit
$p'_{C_y} : X_{k(y)} \to Y_y$ est \uline{birationnel}\hand{, et pour tout point x de $X_{k(y)}$ au dessus}
\hand{d'un pt fermé de X, $p'_{C_y}$ est \uline{étale} en x et \unsure{$\overline{m}$|nr} en les points de $p'^{-1}_{C_y}p'_{C_y}(x)$.}
Signalons la conséquence suivante:

\uline{Corollaire} 14.6. Soit X un schéma algébrique projectif \inserted[typed]{irréductible} \xout{irréductible}
\uline{et séparable}
de dimension n sur un corps infini k. Alors \deleted{X est birationellement} \inserted{$\exists$ un plongement birationnel de X}
\deleted{isomorphe à} \inserted{\deleted{\unsure{dans}} sur} une hypersurface \deleted{dans} P$^{n+1}$ .

On se gardera de croire \deleted{que}, même si X est un sous-préschéma fermé
lisse géom irréd de P, de dimension m-1 = n , \inserted[typed]{que} la projection conique p$_{C_{\hand{y}}}$
soit \inserted[typed]{nécessairement} une \uline{immersion} . En effet, lorsque k est infini, cela impliquerait
qu'il existe une C rationnelle sur k ayant la même propriété, donc que X
est isomorphe à une hypersurface non singulière de \xout{PxxxQxx} P$^{n+1}$. Or
même pour n=1, (\inserted{dans X}courbe algébrique \inserted[typed]{projective} lisse \xout{\unsure{sur}} et connexe sur un corps alg.
clos), il est facile de construire des exemples où X ne peut s'immerger
dans un P$^2$ . De même, dans 14.4. , on se gardera de confondre l'énon-%
cé donné avec l'assertion (en général fausse) que p$_C$ soit lui-même un
monomorphisme (contre-exemple précédent\xout{\unsure{, xx que}}, ou contre-exemple enco-%
re plus évident si X est lisse de dimension m), ou que p$_C$ soit non
ramifié. Pour ce dernier point, on prendra pour s'en convaincre X
sous-\xout{pré}schéma fermé lisse \xout{de dim} irréd de dim m (sur k alg clos disons
si \xout{XxxxxQx} on avait un X $\to$ Q $\simeq$ P$^m$ non ramifié, il serait étale pour
raisons de dimensions, or on peut montrer (cf Chap VIII !) que cela impli-%
que X $\xrightarrow{\sim}$ P$^m$ (P$^m$ étant "simplement connexe). \hand{--} La signification géométri-%
que intuitive de 14.4. est que\xout{x pourxxxvariablexdans} l'ensemble de
ramification de p$_{C_{\hand{y}}}$ est ``variable'' \inserted{sur k}, plus précisément l'ensemble de
ramification de p$_{C_{\hand{y}}}$ pour \hand{y} variable dans \inserted{un ouvert de} Grass$^{m+1}$(\hand{$\bar{k}$}) varie dans X(\hand{$\bar{k}$}), et
n'\uline{admet aucun} \deleted{point fermé de X comme} \uline{``point fixe''} ....

\mspage{90}
